%!TeX program = lualatex
\renewcommand{\thelektion}{L01 \textperiodcentered{} Was ist Informatik?}

\begin{frame}\lessontitle{Lektion 01}{Was ist Informatik?}{Herkunft und Geschichte des Fachs \textperiodcentered{} Skript, Kapitel 1}\end{frame}

\begin{frame}\FDUtitle{Worum geht es in der Informatik?}
Informatik ist weit mehr als das Schreiben von Programmen.
\vspace{3mm}
\begin{center}
\enquote{\textit{Informatik}} = \enquote{Information} + \enquote{Automatik}
\end{center}
\vspace{3mm}
Die Wissenschaft von der \textbf{systematischen} \ldots
\begin{itemize}\setlength\itemsep{4pt}
  \item Darstellung
  \item Speicherung
  \item Übertragung
  \item Verarbeitung
\end{itemize}
\ldots von \textbf{Informationen} mit Hilfe formaler Methoden und technischer Systeme.
\end{frame}

\begin{frame}\FDUtitle{Zwei Ebenen}
\begin{columns}[T,onlytextwidth]
\begin{column}{.46\textwidth}
  \large\bfseries\color{FDUteal}Die Idee
  \vspace{2mm}
  \normalsize\begin{itemize}\setlength\itemsep{6pt}
    \item Algorithmen
    \item Datenstrukturen
    \item Wie löst man ein Problem?
  \end{itemize}
\end{column}
\begin{column}{.46\textwidth}
  \large\bfseries\color{FDUblue}Die Umsetzung
  \vspace{2mm}
  \normalsize\begin{itemize}\setlength\itemsep{6pt}
    \item Hardware, Schaltungen
    \item Netzwerke, Betriebssysteme
    \item Wie baut man die Lösung?
  \end{itemize}
\end{column}
\end{columns}
\vfill\small\color{FDUteal}Dieses Semester begleitet uns beides -- von der Zahl bis zum fertigen Computer.
\end{frame}

{\setbeamertemplate{footline}{\darkfootline}
\begin{frame}\sectionframe[FDUgreen]{Sechs Köpfe, die das Fach geprägt haben}{Was hat Informatik eigentlich mit diesen Menschen zu tun?}
\end{frame}}

\begin{frame}\personcard{Charles Babbage}{1791\,--\,1871, England}{\FDUassets/people/babbage.jpg}{Entwarf die \textbf{Difference Engine} und die \textbf{Analytical Engine} -- mechanische Maschinenideen, die als Vorläufer programmierbarer Computer gelten.}{Als ein Abgeordneter ihn fragte, ob bei falscher Eingabe die richtige Antwort herauskommen könne, entgegnete Babbage nur trocken, er könne \enquote{die Art der Gedankenverwirrung nicht begreifen}, die eine solche Frage hervorrufen könne.}\end{frame}

\begin{frame}\inventionframe{Babbages Difference Engine}{\FDUassets/people/babbage_difference_engine.jpg}{Nachbau der Difference Engine No.\,2 (Science Museum London): eine rein mechanische Rechenmaschine mit tausenden Zahnrädern.}\end{frame}

\begin{frame}\personcard{Ada Lovelace}{1815\,--\,1852, England}{\FDUassets/people/lovelace.jpg}{Schrieb 1843 den ersten veröffentlichten \textbf{Algorithmus} für eine Maschine (Berechnung der Bernoulli-Zahlen) -- gilt daher als erste Programmiererin der Geschichte.}{\enquote{The Analytical Engine has no pretensions whatever to originate anything. It can do whatever we know how to order it to perform.}\\[1mm]$\rightarrow$ eine Maschine tut nur, was man ihr vorgibt -- nicht mehr, nicht weniger.}\end{frame}

\begin{frame}\personcard{Alan Turing}{1912\,--\,1954, England}{\FDUassets/people/turing.jpg}{Beschrieb 1936 mit der \textbf{Turingmaschine} ein universelles Rechenmodell -- die theoretische Grundlage jedes heutigen Computers. Half im 2.\,Weltkrieg, den Enigma-Code zu knacken.}{\enquote{We can only see a short distance ahead, but we can see plenty there that needs to be done.}}\end{frame}

\begin{frame}\inventionframe{Die Enigma}{\FDUassets/people/enigma.jpg}{Eine deutsche Enigma-Verschlüsselungsmaschine: Turings \enquote{Bombe} half, ihren Code zu entschlüsseln -- ein Wendepunkt des Krieges.}\end{frame}

\begin{frame}\personcard{Konrad Zuse}{1910\,--\,1995, Deutschland}{\FDUassets/people/zuse.jpg}{Baute 1941 mit der \textbf{Z3} den ersten funktionsfähigen, frei programmierbaren Computer der Welt -- elektromechanisch, mit Binärzahlen.}{Zuse entwickelte die Z3 in seiner Berliner Wohnung, während in den USA gleichzeitig staatlich finanzierte Grossprojekte wie der ENIAC entstanden.}\end{frame}

\begin{frame}\inventionframe{Die Zuse Z3}{\FDUassets/people/z3.jpg}{Nachbau der Z3 im Deutschen Museum München: rund 2\,600 Relais, gesteuert durch Lochstreifen.}\end{frame}

\begin{frame}\personcard{Claude Shannon}{1916\,--\,2001, USA}{\FDUassets/people/shannon.jpg}{Zeigte 1937, dass sich logische Aussagen (UND, ODER, NICHT) mit elektrischen Schaltkreisen realisieren lassen -- Geburtsstunde der digitalen Schaltungstechnik (siehe Kapitel 4, \enquote{Logische Gatter}).}{\enquote{Information is the resolution of uncertainty.}\\[1mm]Shannon gilt als Begründer der Informationstheorie und prägte den Begriff \textbf{bit}.}\end{frame}

\begin{frame}\personcard{Grace Hopper}{1906\,--\,1992, USA}{\FDUassets/people/hopper.jpg}{Entwickelte den ersten \textbf{Compiler} und machte Programmiersprachen dadurch für Menschen lesbar (Vorläufer von COBOL).}{\enquote{It's easier to ask forgiveness than it is to get permission.}}\end{frame}

\begin{frame}\inventionframe{Der erste \enquote{Bug}}{\FDUassets/people/first_bug.jpg}{9.\,September 1947: Hoppers Team findet eine Motte in einem Relais und klebt sie ins Logbuch -- daher stammt (der Legende nach) der Begriff \textbf{Bug} für einen Programmierfehler.}\end{frame}

{\setbeamertemplate{footline}{\darkfootline}
\begin{frame}\sectionframe[FDUblue]{Woher stammen Zahlen eigentlich?}{Bevor wir zu Computern kommen: wie stellt man überhaupt eine Anzahl dar?}
\end{frame}}

\begin{frame}\FDUtitle{Historischer Hintergrund: Der Ishango-Knochen}
\begin{columns}[c,onlytextwidth]
\begin{column}{.46\textwidth}
\centering
\includegraphics[height=42mm,keepaspectratio]{Grundlagen_Info/13_ZahlendarstellungenUndKodierungen/Figures/ishango-bone}
\end{column}
\begin{column}{.46\textwidth}
Ca.\ 20\,000 Jahre alt, gefunden in der heutigen Demokratischen Republik Kongo.
\vspace{3mm}
\begin{itemize}\setlength\itemsep{5pt}
  \item Deutlich erkennbare, von Menschen eingeritzte Kerben
  \item Vermutete Bedeutung: eine gezählte Anzahl (z.\,B.\ gefangener Fische)
\end{itemize}
\end{column}
\end{columns}
\end{frame}

\begin{frame}\FDUtitle{Von der Kerbe zur Zahl}
\begin{center}
Stellen wir uns vor, die Kerben zählen gefangene Fische.
\end{center}
\begin{equation*}
	\text{drei Kerben} \quad\leftrightarrow\quad \twemoji{tropical fish}\twemoji{tropical fish}\twemoji{tropical fish} \quad\leftrightarrow\quad |||
	\quad\leftrightarrow\quad \text{\textbullet\textbullet\textbullet}
\end{equation*}
\begin{center}
Alle vier Darstellungen verwenden nur \textbf{ein einziges} Symbol $\rightarrow$ \tib{unäre Zahlendarstellung}.
\end{center}
\end{frame}

\begin{frame}\FDUtitle{Unäre Zahlendarstellung: Vor- und Nachteile}
\begin{center}
Was könnten Nachteile dieser Zahlendarstellung sein?
\end{center}
\only<2->{
\vspace{-1mm}
\begin{center}
||||||||||||||||||||||||||||||||||||||
\end{center}
\vspace{-1mm}
\begin{itemize}\setlength\itemsep{3pt}
  \item Wie liest man diese Zahl auf einen Blick?
  \item Die Darstellungslänge wächst \textbf{linear} mit der Zahl -- ungeeignet für grosse Zahlen
  \item Nicht besonders geeignet zum Rechnen
\end{itemize}
}
\end{frame}

\auftragframe{Kapitel 1 (Was ist Informatik?) \& Kapitel 2 (Zahlendarstellungen), Abschnitt \enquote{Ursprünge der Zahlendarstellung}}{%
\aufgabebox{Lesen Sie die Übersicht zur Informatikgeschichte im Skript}\\[3mm]
\aufgabebox{Übungen zur unären Zahlendarstellung}
}


